\documentclass[10pt,a4paper]{amsart}
\usepackage{amsmath,amssymb,amsthm,graphicx}
% Retained preamble for the Stacks Project chapter snapshot in this directory.
% The chapter's own preamble.tex is a copy of the shared Stacks Project preamble; the block
% below is sliced MECHANICALLY, line for line, from it. Nothing was added, removed or
% redefined; the wrapper class, the xr-hyper/hyperref wiring and the externaldocument list
% are omitted because the assembler supplies the document class.
%
% lines taken from preamble.tex: 3, 5, 9, 11, 12, 156, 157, 158, ... (51 lines total)
%
% package coverage restored from the chapter's own preamble (2026-10-13): xy, xy-2cell

\usepackage[all]{xy}
\xyoption{2cell}
\UseAllTwocells

\newenvironment{reference}{\comment}{\endcomment}
\newenvironment{slogan}{\comment}{\endcomment}
\newenvironment{history}{\comment}{\endcomment}
\theoremstyle{plain}
\newtheorem{theorem}[subsection]{Theorem}
\newtheorem{proposition}[subsection]{Proposition}
\newtheorem{lemma}[subsection]{Lemma}
\theoremstyle{definition}
\newtheorem{definition}[subsection]{Definition}
\newtheorem{example}[subsection]{Example}
\newtheorem{exercise}[subsection]{Exercise}
\newtheorem{situation}[subsection]{Situation}
\theoremstyle{remark}
\newtheorem{remark}[subsection]{Remark}
\newtheorem{remarks}[subsection]{Remarks}
\numberwithin{equation}{subsection}
\def\lim{\mathop{\mathrm{lim}}\nolimits}
\def\colim{\mathop{\mathrm{colim}}\nolimits}
\def\Spec{\mathop{\mathrm{Spec}}}
\def\Hom{\mathop{\mathrm{Hom}}\nolimits}
\def\Ext{\mathop{\mathrm{Ext}}\nolimits}
\def\SheafHom{\mathop{\mathcal{H}\!\mathit{om}}\nolimits}
\def\SheafExt{\mathop{\mathcal{E}\!\mathit{xt}}\nolimits}
\def\Sch{\mathit{Sch}}
\def\Mor{\mathop{\mathrm{Mor}}\nolimits}
\def\Ob{\mathop{\mathrm{Ob}}\nolimits}
\def\Sh{\mathop{\mathit{Sh}}\nolimits}
\def\NL{\mathop{N\!L}\nolimits}
\def\CH{\mathop{\mathrm{CH}}\nolimits}
\def\proetale{{pro\text{-}\acute{e}tale}}
\def\etale{{\acute{e}tale}}
\def\QCoh{\mathit{QCoh}}
\def\Ker{\mathop{\mathrm{Ker}}}
\def\Im{\mathop{\mathrm{Im}}}
\def\Coker{\mathop{\mathrm{Coker}}}
\def\Coim{\mathop{\mathrm{Coim}}}
\DeclareMathSymbol{\boxtimes}{\mathbin}{AMSa}{"02}
\def\QCohstack{\mathcal{QC}\!\mathit{oh}}
\def\Cohstack{\mathcal{C}\!\mathit{oh}}
\def\Spacesstack{\mathcal{S}\!\mathit{paces}}
\def\Quotfunctor{\mathrm{Quot}}
\def\Hilbfunctor{\mathrm{Hilb}}
\def\Curvesstack{\mathcal{C}\!\mathit{urves}}
\def\Polarizedstack{\mathcal{P}\!\mathit{olarized}}
\def\Complexesstack{\mathcal{C}\!\mathit{omplexes}}
\def\Pic{\mathop{\mathrm{Pic}}\nolimits}
\def\Picardstack{\mathcal{P}\!\mathit{ic}}
\def\Picardfunctor{\mathrm{Pic}}
\def\Deformationcategory{\mathcal{D}\!\mathit{ef}}
% lane-s11: discard bibliographic \begin{reference} blocks entirely (incl. any \cite inside).
\newsavebox{\limrefbox}
\ifcsname reference\endcsname\expandafter\let\csname reference\endcsname\undefined\fi\ifcsname endreference\endcsname\expandafter\let\csname endreference\endcsname\undefined\fi
\def\citation#1{\relax}
\newenvironment{reference}{\setbox\limrefbox=\vbox\bgroup\let\cite\relax}{\egroup\ignorespaces}
\usepackage{cleveref}
\crefname{originalreference}{original reference}{original references}
\Crefname{originalreference}{Original reference}{Original references}
\makeatletter
\newlabel{algebra-lemma-characterize-dvr}{{00PD}{source}{}{sourcepointer-1}{}}
\newlabel{algebra-lemma-characterize-dvr@cref}{{[originalreference][0][]00PD}{[1][0][]source}{}{}{}}
\makeatother
\crefname{originallemma}{original lemma}{original lemmas}
\Crefname{originallemma}{Original lemma}{Original lemmas}
\makeatletter
\newlabel{lemma-morphism-glueing-near-closed-point}{{0GWT}{source}{}{sourcepointer-2}{}}
\newlabel{lemma-morphism-glueing-near-closed-point@cref}{{[originallemma][0][]0GWT}{[1][0][]source}{}{}{}}
\makeatother
\crefname{originallemma}{original lemma}{original lemmas}
\Crefname{originallemma}{Original lemma}{Original lemmas}
\makeatletter
\newlabel{lemma-Noetherian-dvr-valuative-separation}{{0207}{source}{}{sourcepointer-3}{}}
\newlabel{lemma-Noetherian-dvr-valuative-separation@cref}{{[originallemma][0][]0207}{[1][0][]source}{}{}{}}
\makeatother
\crefname{originalreference}{original reference}{original references}
\Crefname{originalreference}{Original reference}{Original references}
\makeatletter
\newlabel{morphisms-lemma-finite-type-Noetherian-quasi-separated}{{01T7}{source}{}{sourcepointer-4}{}}
\newlabel{morphisms-lemma-finite-type-Noetherian-quasi-separated@cref}{{[originalreference][0][]01T7}{[1][0][]source}{}{}{}}
\makeatother
\crefname{originallemma}{original lemma}{original lemmas}
\Crefname{originallemma}{Original lemma}{Original lemmas}
\makeatletter
\newlabel{lemma-essentially-finite-type-criterion-universally-closed}{{0GWV}{source}{}{sourcepointer-5}{}}
\newlabel{lemma-essentially-finite-type-criterion-universally-closed@cref}{{[originallemma][0][]0GWV}{[1][0][]source}{}{}{}}
\makeatother
\title{Essentially finite type criterion: separatedness part}
\author{The Stacks Project Authors}
\date{stacks-project:limits.tex:local-snapshot-20261012; GNU FDL 1.2 or later}
\begin{document}
\maketitle
\noindent\textit{Editorial note.} No omissions: the diagonal argument is worked; the sibling criteria are delivered cards cited as proved chapter lemmas.\par\medskip
\noindent\textit{Allowed standard prerequisites (editorial).}\begin{itemize}
\item DVR valuative separation. Limits, Lemma lemma-Noetherian-dvr-valuative-separation (same chapter)
\item Ess-ft universally closed criterion. Limits, Lemma lemma-essentially-finite-type-criterion-universally-closed (same chapter)
\end{itemize}\medskip
\ifdefined\hypertarget\hypertarget{sourcepointer-1}{}\fi
\noindent\textit{Source reference.} Original reference 00PD: sources/stacks-chow/tags.txt, source lines 927--927. Cross-chapter standard prerequisite cited by the author inside the counted span; it is used as a pointer and its own proof is not part of this item.\par\medskip
\ifdefined\hypertarget\hypertarget{sourcepointer-2}{}\fi
\noindent\textit{Source reference.} Original lemma 0GWT: sources/stacks-chow/tags.txt, source lines 21083--21083. Cross-chapter standard prerequisite cited by the author inside the counted span; it is used as a pointer and its own proof is not part of this item.\par\medskip
\ifdefined\hypertarget\hypertarget{sourcepointer-3}{}\fi
\noindent\textit{Source reference.} Original lemma 0207: sources/stacks-chow/tags.txt, source lines 2624--2624. Cross-chapter standard prerequisite cited by the author inside the counted span; it is used as a pointer and its own proof is not part of this item.\par\medskip
\ifdefined\hypertarget\hypertarget{sourcepointer-4}{}\fi
\noindent\textit{Source reference.} Original reference 01T7: sources/stacks-chow/tags.txt, source lines 2364--2364. Cross-chapter standard prerequisite cited by the author inside the counted span; it is used as a pointer and its own proof is not part of this item.\par\medskip
\ifdefined\hypertarget\hypertarget{sourcepointer-5}{}\fi
\noindent\textit{Source reference.} Original lemma 0GWV: sources/stacks-chow/tags.txt, source lines 21085--21085. Cross-chapter standard prerequisite cited by the author inside the counted span; it is used as a pointer and its own proof is not part of this item.\par\medskip
\begin{lemma}
\label{lemma-essentially-finite-type-criterion-separation}
Let $S$ be a Nagata scheme (and in particular locally Noetherian).
Let $f : X \to Y$ be a morphism of schemes locally of finite type over $S$.
The following are equivalent
\begin{enumerate}
\item $f$ separated,
\item for any commutative diagram
$$
\xymatrix{
U \ar[r] \ar[d] & X \ar[d]^f \\
C \ar[r] \ar@{..>}[ru] & Y
}
$$
of schemes over $S$ such that
\begin{enumerate}
\item $C$ is a normal integral scheme of finite type over $S$,
\item $U = C \setminus \{c\}$ for some closed point $c \in C$,
\item $A = \mathcal{O}_{C, c}$ has dimension $1$\footnote{It follows
that $A$ is a discrete valuation ring, see
Algebra, Lemma \ref{algebra-lemma-characterize-dvr}.
Moreover, $c$ maps to a finite type point $s \in S$ and
$A$ is essentially of finite type over $\mathcal{O}_{S, s}$.}
\end{enumerate}
then in the commutative diagram
$$
\xymatrix{
\Spec(K) \ar[r] \ar[d] & X \ar[d]^f \\
\Spec(A) \ar[r] \ar@{-->}[ru] & Y
}
$$
where $K = \text{Frac}(A)$ there exists at most one dotted arrow\footnote{By
Lemma \ref{lemma-morphism-glueing-near-closed-point} this is
equivalent to asking there to be at most one dotted arrow
making the first commutative diagram commute.}
making the diagram commute.
\end{enumerate}
\end{lemma}

\begin{proof}
By Lemma \ref{lemma-Noetherian-dvr-valuative-separation} we see that
(1) implies (2). Assume (2). In order to show that $f$ is separated,
we have to show that $\Delta : X \to X \times_Y X$ is closed. By
Morphisms, Lemma \ref{morphisms-lemma-finite-type-Noetherian-quasi-separated}
the morphism $\Delta$ is quasi-compact.
By Lemma \ref{lemma-essentially-finite-type-criterion-universally-closed}
it suffices to show: for any commutative diagram
$$
\xymatrix{
U \ar[rr] \ar[d] & & X \ar[d]^\Delta \\
C \ar[rr]^{(a_1, a_2)} \ar@{..>}[rru] & & X \times_Y X
}
$$
of schemes over $S$ such that
\begin{enumerate}
\item $C$ is a normal integral scheme of finite type over $S$,
\item $U = C \setminus \{c\}$ for some closed point $c \in C$,
\item $A = \mathcal{O}_{C, c}$ has dimension $1$.
\end{enumerate}
then in the commutative diagram
$$
\xymatrix{
\Spec(K) \ar[r] \ar[d] & X \ar[d]^\Delta \\
\Spec(A) \ar[r] \ar@{-->}[ru] & X \times_Y X
}
$$
where $K = \text{Frac}(A)$ there exists some dotted arrow
making the diagram commute. By
Lemma \ref{lemma-morphism-glueing-near-closed-point}
the existence of the dotted arrow in the second diagram
is equivalent to the existence of the dotted arrow
in the first diagram. Moreover, the existence there
is the same as asking $a_1 = a_2$. However $a_1|_U = a_2|_U$, so
by the uniqueness assumption (2) we see that this is true
and the proof is complete.
\end{proof}

\end{document}
